% !TEX root = ../main.tex
\begin{abstract}

Frequent batch auctions (FBAs) are an alternative market structure to continuous limit order books (CLOBs) in which, instead of immediate trade matching, periodic uniform-price clearing occurs over short time intervals. FBAs reduce the value of sub-second latency advantages, which are exploited by high-frequency traders (HFTs). The issue with this market structure is that participants never see the underlying data used to compute the final clearing price; hence, its soundness is questionable. In this paper, we propose a zero-knowledge FBA protocol and provide a construction in which the person announcing the final clearing price at an exchange proves with a custom zk-SNARK—without revealing individual orders—that a published clearing price and allocation satisfy the auction rule. Thus, this protocol preserves participant privacy while providing public verifiability for the end result. We arithmetize the FBA clearing process as vanishing polynomial identities over a multiplicative subgroup, built from reusable gadgets and verified via KZG10 commitments with Fiat-Shamir batching. Beyond theoretical formulation, we implemented the protocol in both Rust/arcworks and Noir/Barretenberg, each proved at 100 to 1,000 price ticks, using both our own synthetic data and real data from Nasdaq TotalView, demonstrating its practical utility.

\end{abstract}